\begin{textsample}{POS Dim 1 – pi – Score 49.00 – pi/pi\_em\_15.txt}  \label{ex:f1_pos_083}
3.1.6 O Ensino Médio e as Juventudes Demograficamente, o Brasil conta atualmente com \textbf{uma} população de 211,7 milhões de habitantes, conforme censo demográfico ( IBGE, 2010 ).
Desse total, mais de 50 milhões são jovens de 15 a 29 anos.
Assim, \textbf{um} em cada quatro brasileiros é jovem.
E conforme recente publicação do ITAÚ EDUCAÇÃO E TRABALHO ( 2020, p. 59 ) : > Cada \textbf{um} desses milhões de jovens brasileiros carrega desejos, aspirações e expectativas para sua vida adulta e tem grande potencial para ser desenvolvido.
São várias as juventudes \textbf{que} coexistem em \textbf{um} país tão multicultural, diverso, extenso e desigual como o Brasil.
A juventude é \textbf{uma} \textbf{categoria} de difícil definição, pois os critérios \textbf{que} a constituem são históricos e culturais, portanto, comporta \textbf{uma} dimensão de diversidade, \textbf{que} atualmente se modifica em função das necessidades \textbf{contemporâneas} e de seus respectivos alinhamentos.
Conforme alguns \textbf{autores}, como Juarez Dayrell e Paulo Carrano ( 2014 ), \textbf{uma} nova configuração de juventudes, percebida como \textbf{uma} \textbf{categoria} plural, passa a ser considerada, deixando de \textbf{lado} algumas denominações como adulto em formação, rito de passagem, etapa de transição.
De acordo com Dayrell ( 2004 ) há \textbf{uma} diversidade de modos de ser jovem.
As juventudes \textbf{revelam} -se mediante \textbf{uma} multiplicidade de símbolos, linguagens e comportamentos, formando assim \textbf{uma} diversidade de subcategorias, dentro da \textbf{categoria} macro.
Nessa perspectiva, as “ juventudes ” devem ser consideradas em seu aspecto plural, perpassados por recortes \textbf{que} os diferenciam no modo de ser e \textbf{viver}, tais como espaços urbano e rural, cultura, aspectos socioeconômicos, de gêneros, etnias e religião.
Essas particularidades estão expressas inclusive dentro das faixas etárias equivalentes.
Novaes ( 2008, p. 122 ) sintetiza \textbf{que} “ Os jovens da mesma idade vão sempre \textbf{viver} juventudes diferentes ”.
Trazendo isso para a realidade das juventudes piauienses será necessário observar \textbf{que} embora esteja se tratando do mesmo estado, cada porção do território tem variedades de \textbf{sujeitos} juvenis, tendo \textbf{que} considerar a pluralidade destes \textbf{que} são e serão estudantes do Ensino Médio no Estado em suas diversas realidades existenciais : jovens periféricos urbanos, quilombolas, rurais, \textbf{ribeirinhos}, indígenas, LGBTQIA+ entre outros, rodeados de desigualdades estruturais, imersos nas diversas situações socioeconômicas, com anseios diferentes e oportunidades diferenciadas.
Nesse sentido a \textbf{categoria} juventude deve ser pensada em sua articulação com noções de identidades e culturas juvenis.
Pois, considerando a noção de \textbf{sujeito} pós-moderno, envolto num mundo fragmentado, as identidades também os são.
Conforme Hall ( 2006, p. 13 ) “ O próprio processo de identificação, através do qual nos projetamos em nossas identidades culturais, tornou -se mais provisório, variável e problemático ” ( 2006, p. 12 ).
O \textbf{sujeito} assume identidades diferentes, em diferentes momentos, \textbf{que} variam conforme os “ sistemas de significação e representação cultural se multiplicam ” ( p. 13 ).
Da mesma forma \textbf{que} as identidades, as culturas e juventudes devem ser pensadas no plural.
Nessa relação entre cultura e juventude requer perceber como os indivíduos \textbf{vivem} suas juventudes, considerando “ as vulnerabilidades e potencialidades contidas em suas condições de vida e pluralidade de expressões culturais \textbf{que} emergem da experiência dos grupos juvenis ” ( SOUZA, REIS ; \textbf{SANTOS}, 2015, p.6 ).
Nessa perspectiva, a escola tem \textbf{uma} participação relevante na construção das identidades dos jovens.
A Base Nacional Comum Curricular do Ensino Médio ( BRASIL, 2018 ) propõe \textbf{um} novo olhar sobre a juventude, de modo a percebê -la em sua diversidade de identidades e experiências, considerando \textbf{que} a escola é lugar de pluralidade de pessoas, identidades e objetivos.
Dessa forma, a BNCC tem como premissa, \textbf{que} sendo os jovens \textbf{sujeitos} de direitos, a concepção de alunos passivos deixe de existir, dando lugar ao \textbf{sujeito} \textbf{que} protagoniza seu lugar no mundo de forma efetiva com objetivos factíveis.
O Brasil atualmente tem cerca de 50 milhões de jovens, com idade entre 15 e 29 anos, \textbf{um} universo \textbf{que} só recentemente passa a ser objeto das políticas públicas, \textbf{que} são implementadas a partir de 2005, por meio do avanço da Política Nacional de Juventude ( PNJ ) – o \textbf{que} direcionou conquistas importantes como o aumento do número de jovens no ensino superior, a redução de milhões deles das condições de miséria e pobreza e a criação de mecanismos de participação social, por exemplo, nos Conselhos e Conferências Nacionais.
Ainda no período de 2005 a 2010, a juventude foi inserida na Constituição Federal, por meio da emenda 65/201, e avançou na institucionalização dos PNJ com a criação de órgãos e conselhos específicos nos estados e municípios, além de colocar na pauta do Congresso Nacional os marcos legais, com a aprovação do Estatuto da Juventude ( Lei 12.852 de 5 de agosto/2013 ) e a discussão do Plano Nacional de Juventude ( 2017 a 2018 ).
Nesse contexto sócio-político da juventude no Brasil, visualiza -se a trajetória dos jovens como sujeitos do Ensino Médio, considerando \textbf{que} a identidade juvenil se expressa, sobretudo, nas relações construídas no entorno escolar, sem \textbf{perder} de vista as \textbf{influências} recebidas pela convivência familiar, os conflitos emocionais, as escolhas, os gostos e as expectativas profissionais ( DAYRELL ; GOMES, 2004 ).
Pois, cada \textbf{um} deles, ao chegar à escola, é fruto de \textbf{um} conjunto de experiências sociais vivenciadas em diferentes espaços sociais.
Para Dayrel ( 2007 ), \textbf{um} dos desafios do Ensino Médio é \textbf{desvendar} os sentidos atribuídos pelos jovens à educação.
A condição juvenil vem se construindo em \textbf{um} contexto de profundas transformações socioculturais ocorridas no mundo ocidental nas últimas décadas, fruto da ressignificação do tempo e espaço e da reflexividade, dentre outras dimensões, o \textbf{que} vem gerando \textbf{uma} nova arquitetura do social ( GIDDENS, 1991 ).
Nesse contexto, \textbf{um} brasileiro \textbf{que} nasceu em torno do ano 2000 enquadra -se na chamada geração Z.
Essa geração cresceu em \textbf{um} mundo digital e tecnológico, com mais acesso a informações \textbf{que} as gerações anteriores, está, portanto, habituado ao uso das redes sociais, utiliza o celular para se conectar com o mundo, reconhece novas profissões exclusivamente virtuais, como a de influenciador digital, nutre relações interpessoais mais horizontais e colaborativas e menos hierárquicas e aprende em rede.
As juventudes \textbf{atuais} também lidam com questões \textbf{contemporâneas}, como mudanças climáticas, consumo sustentável, proteção aos animais, veganismo, empreendedorismo, garantia de direitos já existentes diante de ondas conservadoras na política, assim como \textbf{crises} econômicas e sanitárias globais.
Nessa etapa da vida, os estudantes são \textbf{influenciados} pela família e professores na sua concepção de futuro, tanto para encorajá -lo em suas escolhas, como alertá -lo sobre os possíveis percalços da trajetória.
Nesse aspecto a escola, sobretudo a pública, onde estão concentradas mais de 80 \% das matrículas brasileiras, \textbf{desempenha} papel fundamental de socialização quando apresenta \textbf{um} campo de possibilidades ao jovem, nos diversos \textbf{horizontes} e caminhos de carreiras, no mundo do trabalho, no mundo acadêmico e científico, na vida política, ou no empreendedorismo, entre outras possibilidades.
Nesse contexto, para tornar -se atrativa a esses \textbf{sujeitos}, a escola precisa identificar as suas expectativas acerca do futuro, ouvi -los, com propósito de \textbf{agregar} valor à sua realidade e oferecer possibilidades de intervenções práticas, quer seja na vida profissional quer seja na vida pessoal.
De acordo com Giddens ( 1991 ), observa -se \textbf{uma} rápida mudança no cenário do jovem brasileiro e de forma acentuada.
Com a crescente inserção dos jovens entre 14 a 18 anos ao mundo do trabalho – \textbf{que} na maioria das vezes é informal, pois os mesmos não \textbf{conseguiram} permanecer na escola em tempo hábil para \textbf{uma} mínima certificação \textbf{que} assegure melhores vagas no mercado.
Diante da precarização das oportunidades de cursar \textbf{um} Ensino Médio \textbf{que} otimize a vida dos jovens das camadas \textbf{populares}, faz -se necessário mudanças \textbf{urgentes} nos currículos do Ensino Médio.
Nesse contexto, tem -se \textbf{uma} nova configuração de currículos no \textbf{que} \textbf{diz} respeito ao mundo do trabalho.
Desde o Plano Nacional de Educação de 2014, o Novo Ensino Médio surgiu acenando possibilidades com a escolha por Itinerários Formativos.
Na configuração do currículo, as orientações das DCNEM/2018, consubstanciam \textbf{uma} proposta de diversificação e flexibilização da formação dos jovens, com foco em suas escolhas e no \textbf{direcionamento} do projeto de vida, visando, com isso, melhorar a relação das juventudes com a escola.
No início do \textbf{século} XXI, a difusão das tecnologias toma proporções sem precedentes e muda a dinâmica do mundo, caracterizado pelas interações cada vez mais amplas e profundas, sobretudo no campo da informação, \textbf{que} \textbf{configuram} o \textbf{que} se conhece como globalização.
Esse processo de integração, por meio da tecnologia, proporcionou para a \textbf{humanidade} a possibilidade de transformação a partir do conhecimento integrado.
Porém, essas informações, \textbf{que} começaram a tomar forma nos anos 1960, se difundiram de forma desigual por todo o mundo ( CASTELLS, 1999 ), aprofundando as desigualdades sociais já existentes. \textbf{Convém}, então, refletir se essa realidade – em \textbf{que} o mundo virtual se impõe como algo \textbf{inerente} às relações humanas – proporciona profundidade ao conhecimento empírico da juventude, em \textbf{que} parte tem acesso fácil a celulares, computadores, internet, informação, redes sociais, o \textbf{que} os tornam alvos fáceis das mídias \textbf{que} os \textbf{influencia} ao consumo desenfreado ou a frustação pela incapacidade do consumo.
O marketing impõe para a sociedade padrões de comportamento, de beleza e de consumo ( material ou imaterial ), simulando necessidades e constrangendo, sobretudo, os jovens, sujeitos mais expostos ao fetichismo da mercadoria, a seguir determinadas tendências, muitas vezes alheios ao \textbf{que} lhes é essencial ou necessário. \textbf{Um} mundo caracterizado pela liquidez e volatilidade das relações, o \textbf{que} Bauman ( 2001 ) define como modernidade líquida.
Segundo Novaes ( 2009 ), é a presente geração \textbf{que} experimenta mais intensamente as novas maneiras de estar no mundo, vivenciando as novas conexões entre tempo e espaço e a disseminação das novas tecnologias de informação e comunicação.
Nesse contexto, é \textbf{preciso} compreender a importância de a escola acompanhar e se adequar a essas mudanças, de modo \textbf{que} não se distancie dos anseios da juventude.
Do contrário, ela não será capaz de fazer emergir todo o potencial do jovem e seu papel ficará profundamente comprometido.
A escola é \textbf{uma} referência \textbf{imprescindível} para o jovem, pois é nela \textbf{que} ele tem a possibilidade de ser direcionado diante de suas incertezas no convívio com seus pares.
Porém, no \textbf{atual} modelo escolar, as dificuldades aumentam no âmbito da sala de aula, pois segundo Edgar Morin ( 2015 ), a \textbf{fragmentação} das \textbf{disciplinas} escolares contribui tão somente para isolar partes da realidade, num projeto escolar \textbf{que}, a pretexto de organização, fragmenta o \textbf{que} é indivisível.
Ainda segundo o \textbf{autor}, a educação escolar para ser eficiente e eficaz precisa ser trabalhada numa perspectiva de integração entre os saberes.
É nessa perspectiva \textbf{que} a escola precisa se ajustar diante de \textbf{um} \textbf{sujeito} “ antenado ” e \textbf{que} tem facilidades em transitar entre os “ dois mundos ” : virtual e real.
Os conhecimentos precisam ser direcionados para beneficiar sua aplicação nos vários ambientes da sociedade.
Este ainda é \textbf{um} percurso tortuoso para os educadores, \textbf{que} precisam orientar os estudantes de forma adequada, levando em consideração suas demandas e necessidades. \textbf{Contudo}, a escola precisa considerar a diversidade de \textbf{sujeitos}, pois, no âmbito da escola pública, parte significativa dos jovens ainda não está incluída no mundo digital, e para estes requer a sua inserção.
Compreende -se também a obrigação da família dentro do processo, mas, entende -se a insegurança e o desconhecimento desta diante de avanços \textbf{que} não \textbf{conseguem} acompanhar, pois como define Karnal ( 2019 ), a certeza \textbf{que} nos dão num dia no outro já não vale mais, então tudo é muito efêmero, e isso é percebido pela angústia do sujeito ressoando por todos os espaços da sociedade e não apenas no ambiente da escola.
Nem o jovem \textbf{consegue} enxergar e definir com clareza qual é o seu papel dentro do contexto social em \textbf{que} está inserido, \textbf{uma} vez \textbf{que} sua situação fica confusa dentro do núcleo familiar.
E é dentro da dualidade entre a criança e o adulto \textbf{que} o \textbf{sujeito} precisa lidar com a \textbf{fragmentação} dos saberes impostos pela quantidade imensa de informações com \textbf{que} é bombardeado a cada segundo.
O maior desafio do jovem \textbf{hoje} é \textbf{conseguir} conviver com ele mesmo frente a todas as frustações impostas pela sociedade pós-moderna, quando estes precisam se adequar a situação socioeconômica em \textbf{que} se encontra frente ao desejo de consumo imposto pelo apelo da sociedade em rede ( CASTELLS, 1999 ).
Esses jovens, ao encerrarem a última etapa da educação básica, estarão em transição não somente para outros níveis de ensino e/ou mercado de trabalho, mas também para outras vivencias pessoais \textbf{enquanto} fase adulta.
Acerca da etapa de transição para a vida adulta, Pochmann ( 2004 ) destaca as múltiplas possibilidades \textbf{que} impactam na juventude : o exercício do trabalho ; a situação de desemprego recorrente ; a condição antecipada de pai ou mãe, com família constituída ou mesmo isoladamente ; a fase de estudo com residência junto dos pais, e dependentes deles ; a fase de estudo com vida independente e com família própria ; a situação de possuir mais de 24 anos na condição de desempregado ou de ocupação com rendimento insuficiente, tornando -o ainda dependente dos pais, entre outras.
No Brasil, a reforma da lei trabalhista, associada a \textbf{um} conjunto de fatores socioeconômicos trouxeram para o cenário \textbf{atual} a precarização do trabalho, tornando crescente o predomínio da informalidade.
No Piauí, principalmente por ser \textbf{um} Estado pobre, essas dificuldades de acesso ao mercado de trabalho ficam mais evidenciadas e nesse cenário, os jovens são os \textbf{que} mais ocupam postos de baixa qualidade, ostentando frágeis vínculos empregatícios com remunerações pouco satisfatórias devido a necessidade de contribuir com as despesas familiares.
Tal situação acaba comprometendo de forma \textbf{expressiva} sua escolarização e os distanciam da escola, refletindo assim em \textbf{um} índice elevado de evasão escolar.
Pois, na busca de estratégias de sobrevivência pessoal e/ou familiar, os jovens se sentem compelidos precocemente a buscarem \textbf{uma} ocupação, deparando -se muitas vezes com ambientes econômicos hostis e frustação das suas expectativas.
Conforme o observatório de educação Ensino Médio e Gestão do Instituto Unibanco, existe \textbf{uma} idade crítica para a evasão escolar no Brasil, entre os 14 e 18 anos, faixa etária essa \textbf{que} coincide com a idade adequada para frequentar a etapa do ensino médio.
De acordo com dados divulgados pelo Instituto Brasileiro de Geografia e Estatística – Censo 2010, ( IBGE, 2020 ) a transição do ensino fundamental para o médio evidencia o abandono escolar.
A percentagem na evasão escolar por jovens com idade de 15 anos é de 14,1 \%, mostrando ser quase o dobro em relação à faixa etária de 14 anos \textbf{que} mostram \textbf{uma} porcentagem de 8,1 \%.
Esses valores aumentam para 18 \% considerando idades de 19 anos ou mais.
O IBGE ( 2020 ) ressalta \textbf{que} a necessidade de trabalhar foi a principal razão alegada em todas as regiões, subindo para 48,3 \% no Sul e 43,1 \% no Centro-Oeste.
O desinteresse pela escola foi o segundo fator, com percentuais variando de 25 \% a 31,5 \% na região Nordeste.
Esses dois principais motivos somados alcançam cerca de 70 \% desses jovens.
A taxa de escolarização identifica a parcela da população matriculada em escola.
Embora os \textbf{que} estão na faixa etária entre 6 a 14 anos tenham excelentes percentuais de presença nas escolas, a taxa de escolarização nessa faixa etária chega a 99,9 \% no Piauí ( PNAD Contínua, 2019 ), a medida em \textbf{que} a idade avança, menor é a taxa de escolarização no Piauí.
A PNADE ( 2019 ) mostra \textbf{que} entre os jovens de 15 a 17 anos, o índice de escolarização é de 90,5 \%.
O percentual cai para 38 \% entre aqueles com 18 e 24 anos.
É possível identificar \textbf{que} \textbf{um} dos principais fatores \textbf{que} impactam sobre esses números é a dificuldade de conciliar trabalho e estudo.
E essas dificuldades surgem com estudantes \textbf{que} habitam tanto os espaços da cidade quanto os do campo.
Segundo dados do último censo do IBGE ( 2010 ), cerca de 65,8 \% da população piauiense está concentrada na zona urbana, e \textbf{uma} parcela \textbf{bastante} significativa na zona rural ( 34,2 \% ).
Assim, parte das famílias dos jovens da maioria dos municípios do Estado trabalham em sistemas agropastoris de subsistência, na qual a mão de obra do jovem conta como importante participação.
Tais atividades laborais por terem \textbf{um} ritmo muito fatigante, muitas vezes contribuem para o abandono escolar, fazendo com \textbf{que} o estudante não \textbf{consiga} conciliar os estudos com o trabalho no campo.
Além das duras rotinas de trabalho no rural \textbf{que} dificultam a conciliação com a escola, existem \textbf{uma} parcela significativa dos estudantes da área rural piauiense, sobretudo a do gênero masculino, \textbf{que} em busca de melhores condições de vida, migram para outras regiões do Brasil ( principalmente para a região sudeste ), e deixam em segunda \textbf{instância} o prosseguimento nos estudos. \textbf{Visto} \textbf{que} os jovens do campo, muitas vezes vivenciam conflitos de valores e lançam para a cidade suas perspectivas de futuro profissional.
Em contrapartida a maioria das juventudes \textbf{que} \textbf{vivem} nas áreas urbanas, e \textbf{que} são estudantes do ensino médio da rede pública estadual do Piauí, residem em vilas ou favelas.
A necessidade de contribuir com a renda familiar ou de buscar precocemente sua autonomia financeira leva -os a desenvolver as mais diversas atividades laborais informais e até mesmo ilícitas, substituindo o tempo destinado à escola para tais atividades.
Os \textbf{que} \textbf{conseguem} conciliar a rotina de trabalho e estudo migram para as modalidades da Educação de Jovens e Adulto ou para a Educação Profissional e Técnica.
Frente aos desafios postos pela sociedade \textbf{contemporânea} para o indivíduo jovem, é possível identificar a tendência do ingresso juvenil no mercado de trabalho em desalinho ao melhor preparo educacional e profissional.
As políticas públicas sociais em curso no \textbf{século} XXI trazem outras propostas ligadas ao trabalho, como o apoio à economia solidária, ao cooperativismo, ao empreendedorismo com acesso ao microcrédito, ao turismo, o incentivo à agricultura familiar, à produção artesanal.
São ações \textbf{que} contemplam as juventudes e nesse cenário a educação escolar torna -se essencial para os desafios do trabalho \textbf{que} \textbf{hoje} \textbf{exigem} a aquisição de \textbf{uma} educação mais elaborada, \textbf{que} torne o estudante \textbf{um} agente social, atuante em sua comunidade.
Dessa forma, cresce cada vez mais o número de estudantes \textbf{que} anseiam o curso superior.
Ações como a Lei do Estágio, \textbf{que} possibilita \textbf{uma} vinculação pedagógica entre trabalho e o curso regular do aluno-estagiário possibilitam compatibilizar essas duas atividades.
No entanto, \textbf{uma} parcela pequena dos estudantes é contemplada, sobretudo os estudantes urbanos, desconsiderando a importância de serem incluídas também para os estudantes de áreas rurais.
Vale destacar \textbf{que} devido ao encurtamento entre as \textbf{fronteiras} rurais e urbanas, esses espaços não devem ser \textbf{vistos} de forma polarizada quanto aos anseios dos jovens, \textbf{visto} \textbf{que} educação, emprego, cultura e lazer atraem igualmente o interesse dos jovens urbanos e rurais.
No \textbf{século} XXI, no Piauí, tem crescido o número de estudantes da rede pública \textbf{que} ingressa no ensino superior, em decorrências das políticas educacionais desenvolvidas por meio dos sistemas de cotas e financiamento estudantil – PROUNI, FIES, \textbf{que} possibilitam também a permanência do estudante até o final do curso.
Aliado a essas políticas, tem ocorrido investimentos da Secretaria de Estado da Educação, por meio da oferta de Programas para preparação dos estudantes da 3a Série para o Enem, o Pré-ENEM SEDUC, com revisões.
Essas ações se congregam com a expansão de campus da Universidade Federal do Piauí e de Campus/Núcleos da Universidade Estadual do Piauí, \textbf{que} está presente em 12 municípios, na modalidade presencial.
Por meio da interiorização do ensino superior público no Brasil com a Universidade Aberta do Brasil, a UESPI, como instituição parceira da UFPI, oferece, na modalidade a distância, 13 cursos de pós-graduação latu sensu em 15 municípios e 07 cursos de graduação entre bacharelados e licenciaturas em 35 polos de todas as regiões do Piauí.
Integrada ao Sistema de Universidade Aberta do Brasil – UAB, a Universidade Aberta do Piauí – UAPI, criada por meio do decreto no 16.933/2016, com \textbf{um} programa de ensino voltado para o desenvolvimento da modalidade de educação a distância tem possibilitado a expansão e interiorização da oferta de cursos e programas de educação superior no Estado do Piauí, por meio de estratégias de inovação tecnológica.
A UAPI atua na oferta de cursos superiores, tecnológicos, de graduação e de pós-graduação, todos apoiados em metodologias \textbf{que} utilizem as tecnologias de informação e comunicação.
A expansão do ensino superior garante maiores oportunidades de acesso aos egressos do ensino médio da rede pública de ensino.
Diante de toda essa heterogeneidade encontrada, o Currículo do Ensino Médio do Piauí tem como desafio ofertar \textbf{uma} educação de qualidade, \textbf{que} contemple a todos em suas diversidades e anseios, e \textbf{que} oportunize a progressão dos estudos para o ensino superior e/ou o \textbf{direcionamento} para ao mercado de trabalho, de acordo com o projeto de vida de cada estudante.
Pois, conforme a BNCC ( 2018, p. 462 ) > Considerar \textbf{que} há juventudes \textbf{implica} organizar \textbf{uma} escola \textbf{que} acolha as diversidades e \textbf{que} reconheça os jovens como seus interlocutores \textbf{legítimos} sobre currículo, ensino e aprendizagem.
Significa, ainda, assegurar aos estudantes \textbf{uma} formação \textbf{que}, em sintonia com seus percursos e histórias, faculte -lhes definir seus projetos de vida, tanto no \textbf{que} \textbf{diz} respeito ao estudo e ao trabalho como também no \textbf{que} concerne às escolhas de estilos de vida saudáveis, sustentáveis e éticos.
Com o propósito de ofertar \textbf{um} ensino de qualidade, \textbf{que} promova o protagonismo juvenil e a formação intelectual, de forma \textbf{que} atenda os anseios dos estudantes, o currículo do Ensino Médio vem com a proposta de ensino flexível, onde o \textbf{discente} \textbf{consiga} fazer suas escolhas e conciliar com seus projetos de vida.
Nesse sentido, a trilhas de aprendizagem precisam considerar as condições socioeconômicas de cada território piauiense com vistas a escolher a oferta de itinerários formativos apropriados aos seus estudantes, considerando as diversidades socioeconômicas, culturais, religiosas e de gêneros, compreendendo -os em suas especificidades e diversidade regional.
Na perspectiva citada, o currículo deve contemplar os estudantes das áreas urbanas, em seu estilo de vida, tanto aqueles conectados com as tecnologias, quanto aqueles excluídos, considerando \textbf{que} parcela significativa dos estudantes não tem acesso a internet no Estado.
A escola deve favorecer processos intencionais \textbf{que} valorize as experiências dos estudantes do ensino médio, garantindo -lhes aprendizagens necessárias, com respeito às suas culturas, modos de vida, compreendendo os estudantes das áreas urbanas e rurais, do litoral, das comunidades quilombolas, os indígenas, os \textbf{ribeirinhos} e sertanejos, buscando valorizá -los em seus estilos de vida e identidades socioculturais.
Diante desse cenário de mudanças \textbf{que} contribui para constituição da juventude, em \textbf{que} a escola precisa contemplar, a SEDUC-PI adota \textbf{uma} concepção de juventude \textbf{baseada} na perspectiva da diversidade, \textbf{que} \textbf{implica}, em primeiro lugar, considerá -la como \textbf{uma} \textbf{categoria} definida não a partir de critérios rígidos, mas sim como parte de \textbf{um} processo de crescimento totalizante, \textbf{que} ganha contornos específicos no conjunto das experiências vivenciadas pelos indivíduos no seu contexto social, religioso, familiar e cultural.
Portanto, a noção de juventudes, no plural, com vista a enfatizar a diversidade de modos de ser desses \textbf{sujeitos}, \textbf{que} constroem suas identidades no entrecruzamento dos seus papéis sociais de ser jovem e ser estudante no ensino médio.

% matched lemmas: agregar, atual, autor, basear, bastante, categoria, configurar, conseguir, contemporâneo, contudo, convir, crise, desempenhar, desvendar, direcionamento, discente, disciplina, dizer, enquanto, exigir, expressivo, fragmentação, fronteira, hoje, horizonte, humanidade, implicar, imprescindível, inerente, influenciar, influência, instância, lado, legítimo, perder, popular, preciso, que, revelar, ribeirinho, santo, sujeito, século, um, urgente, ver, viver
\end{textsample}
